\section{Introduction}\label{sec:introduction}

Let \(f\in\mathbb C\{z_1,\ldots,z_N\}\) be a nonzero convergent
power series with \(f(0)=0\). The hypersurface germ \(V(f)\) is
\emph{reduced} when \(f\) has no repeated irreducible factor. Its
\emph{multiplicity} is
\(\operatorname{ord}_0f\), the least total degree of a nonzero term
of \(f\). Two hypersurface germs are \emph{ambiently homeomorphic}
if a homeomorphism between neighborhoods of the origin, fixing the
origin, takes one zero set onto the other. Equivalently, their ambient
pair germs \((\mathbb C^N,V(f),0)\) are homeomorphic.

Zariski's multiplicity question asks whether this embedded
topological germ determines multiplicity. He posed the question
explicitly in 1971 \cite{Zariski1971}, following his work on the
topology and equisingularity of plane curves
\cite{Zariski1932,Zariski1965}. In that classical setting, the
embedded topology determines the characteristic data of the branches
and their mutual intersection numbers, and hence the multiplicity;
see also the proof of \cite[Theorem~3.1]{KoikeParusinski2010}.
The higher-dimensional question asks whether this relation between
an equation's lowest degree and its embedded topology persists.

For an isolated critical point, Milnor's fibration gives a way to
study that topology through the \emph{link}
\(V(f)\cap S^{2N-1}_\varepsilon\) on a sufficiently small sphere
and the associated open book, whose complement fibration is
\(f/|f|\) \cite{Milnor1968}.
We answer Zariski's question negatively in four complex variables.

\begin{theorem}\label{thm:main}
There are reduced convergent holomorphic function germs
\(f_4,f_5\colon(\mathbb C^4,0)\to(\mathbb C,0)\), each with an isolated
critical point at the origin, such that
\[
   \operatorname{ord}_0f_4=4,\qquad
   \operatorname{ord}_0f_5=5,
\]
and a homeomorphism germ of ambient pairs
\[
   (\mathbb C^4,V(f_4),0)
      \;\cong\;
   (\mathbb C^4,V(f_5),0).
\]
\end{theorem}

The Milnor open book has a compact page whose middle-dimensional
cycles carry an integral Seifert form: it records the linking of a
cycle on a reference page with a cycle transported to a positive
neighboring page. For the open books on \(S^7\) used here, a
classification theorem turns a congruence of these integral forms
into an equivalence of the embedded links. The proof is therefore
organized around obtaining that integral congruence while retaining
the different orders of the equations.

\subsection*{Context and predecessors}

A major line of work concerns multiplicity within a deformation.
The Milnor number \(\mu\) is the dimension of the local algebra
defined by the first derivatives; for a convergent isolated germ it
is also the rank of the middle homology of the Milnor page.
Greuel and O'Shea independently proved equimultiplicity for
\(\mu\)-constant deformations of an isolated quasihomogeneous
singularity \cite{Greuel1986,OShea1987}.
Fernández de Bobadilla and Pełka proved the general family theorem:
every coefficientwise continuous path of formal power series with
finite constant Milnor number has constant multiplicity
\cite[Theorem~1.1]{BobadillaPelka2024}. The germs of
Theorem~\ref{thm:main} therefore cannot be endpoints of such a path.

The distinction between a family and an arbitrary pair is substantial.
Fernández de Bobadilla and Pełka's Example~8.6 gives a four-variable comparison:
double suspension of Du Bois--Michel plane curves with isometric
integral Seifert forms yields topologically equivalent germs of
equal multiplicity two, with no connecting \(\mu\)-constant family
\cite[Example~8.6]{BobadillaPelka2024}. Du Bois and Michel's plane
curves have different embedded types despite their isometric forms
\cite{DuBoisMichel1994}.

Another positive result retains geometric information beyond the
smooth embedded link. McLean proved that two isolated polynomial
singularities with embedded contactomorphic links have equal
multiplicity \cite[Theorem~1.1]{McLean2019}: the ambient sphere
map must preserve the natural contact structures as well as the
links. His Floer-theoretic characterization of multiplicity is a
starting point for the family theorem of Fernández de Bobadilla
and Pełka. The comparison in this paper uses smooth open books and
their integral Seifert forms.

Koike and Parusiński proposed high-power suspension of plane curves
with isometric Seifert forms as a route to topologically equivalent
four-variable germs, asking whether their Fukui numerical sets---the
sets of orders attained along analytic arcs---could differ. Their test of the Du Bois--Michel family gives
equal sets \cite[Section~5, pp.~294--296]{KoikeParusinski2010}. The
smooth sum and tensor calculation underlying our construction belongs
to the product theory of Sakamoto and Kauffman--Neumann
\cite{Sakamoto1974,KauffmanNeumann1977}.

A separate article by OpenAI, \emph{Ambiently homeomorphic isolated
hypersurfaces of multiplicities two and three}, constructs
real-coefficient, positive-weight homogeneous polynomials in one
finite complex dimension \(N>3\), divisible by eight, with orders
two and three and homeomorphic ambient germs
\cite[Theorem~1.1]{OpenAI2026Multiplicity23}. That construction uses weighted
homogeneous sums and spherical-link classification. The present
four-variable construction uses convergent equations in analytic
charts, a finite-map lattice calculation, and the nonspherical
open-book theorem. The two arguments are independent, and their
dimension, equation format, and exact multiplicities are distinct.
The ambient three-variable case remains untreated by either result.

\subsection*{Proof strategy}

The construction starts from the plane curve
\[
 h(X,T)=T(X-T^2)(X-2T^2)
\]
and two degree-four finite maps \(\pi_0,\pi_1\) obtained by projecting
smooth graph surfaces to the \((X,T)\)-plane. Their pullbacks \(d_i=h\circ\pi_i\)
are six-branch plane curves of the same embedded type, so their
integral plane Seifert forms are isometric. Thus the plane curves
themselves already have the same topology, whereas Du Bois and Michel's
examples separate topology from the abstract Seifert form. Here the
distinction lies in the maps to a common base curve, whose branch-degree
patterns differ. We retain their induced homology maps separately from
the abstract plane-form isometry. Section~\ref{sec:plane-cover} proves
that isometry and the separate pushforward and transfer identities.

A small rescaling of each graph equation smooths the cone \(yz=w^4\).
On each resulting smooth fourfold we use the function
\(h+w(y^a+z^b)\), with one common choice of \(a,b\). In
Section~\ref{sec:ambient-germs}, that choice is fixed before use;
analytic elimination gives orders five and four, and the subsequent
arc argument proves isolation. To compute the Milnor pages, we cut
each into an inner smoothing region and an exterior inherited from
the cone. An exterior cycle closes with an inner plane cycle only
when their boundary residues agree modulo four. This condition
depends on the pushforward for the particular graph map and is not
determined by the abstract plane Seifert form.
Sections~\ref{sec:tube}--\ref{sec:filling-summands} establish the
condition by a signed Mayer--Vietoris cut and an oriented Bockstein
calculation, then compute the two filling pairings in one convention.
The quotient by the cone's cyclic degree-four cover and its signed
integral gluing are established in this tube calculation. The resulting
lattice and form belong to the actual local Milnor page, as required
for the link theorem.

Two boundary directions of the cone page then give different
splittings of this integral lattice. Their coprime indices make
the associated rational maps integral at complementary sets of
primes, including the dyadic prime. Section~\ref{sec:fourfold-arithmetic}
compares those maps on the relevant quadratic and Hermitian
eigenspaces, lifts their discrepancy to the spin and special
unitary groups, and checks the real and lattice-stabilizer hypotheses
for strong approximation
\cite{Platonov1969,Platonov1970}
\cite[Theorem~2.3]{Rapinchuk2014}. This produces one integral
Seifert-form congruence.

Finally, Kato's classification of simple spinnable structures on
\(S^7\) applies to these Milnor open books without requiring their
bindings to be homotopy spheres \cite[Theorem~B]{Kato1974}.
Kato credits Durfee's independent work \cite{Durfee1974};
Saeki explains the singularity-facing formulation
\cite[Definition~2.1 and pp.~698--699]{Saeki1999}. After the page and
orientation hypotheses are checked, the integral congruence gives
an equivalence of embedded links. The analytic ambient-pair cone theorem
\cite[Theorem~1.4]{Massey2015LeCycles} extends it to the ambient pair
germs; this is the analytic form of the conical construction used by
Milnor \cite[Section~2]{Milnor1968}.
